\section{Limits and non-claims}

The paper does not claim that G\"odel's incompleteness theorems are false. It does not claim that every formal calculus, every reasoning practice, or every metamathematical coding scheme satisfies the target regime. It does not deny external metamathematics, explicit scope change, or undecidability phenomena that live outside same-domain inferential governance. It does not show that no broader completion project is possible after enlarging scope, changing the admissibility regime, or importing new standing-bearing structure. It also does not treat bare derivation-trace systems, coding layers, or external proof checkers as automatically belonging to the target. They enter the argument only if their same-domain readback can alter inferential life on the fixed domain.

Its claim is narrower and stronger inside that narrowness. Once the fixed domain is already forced into the coherent-reasoning kernel, all same-domain governance-relevant closure threats must survive into standing classification or coherent legal trajectory structure. The coherent reasoning layer shows that no hidden path-level third locus remains. The code-role theorem shows that governance-relevant coding is hidden gate work. The route-exhaustion theorems show that no same-domain G\"odel threat architecture survives those constraints.


\paragraph{Provisional practice and repaired proofs.} The framework does not deny exploratory mathematics, seminar sketches, blackboard heuristics, unpublished drafts, or provisional codings. Before they are taken as standing-bearing same-domain acts, they are provisional practice or skin-level bookkeeping. Once a sketch, draft proof, or coding is entered as a standing-bearing same-domain act, however, the kernel applies immediately. A later corrected proof is therefore a fresh act, or an explicit reset of scope or target, rather than a same-act repair inside the retained lineage. This is exactly the sense in which the theory classifies repaired proofs without forbidding mathematical practice.

\paragraph{Formalization status.} The prose theorem chain is primary. A companion Lean4 development serves as evidential audit support for the fixed-kernel consequence layer used here, but detailed theorem-ledger tables and provenance crosswalks are moved out of the primary manuscript package and retained in the audit companion. The body therefore asks to be read as mathematics first and formalization-supported second.